% SOURCE: Chinese PDF viewer pp. 86 (bottom)--91.
\section{Coupling of Two Angular Momenta}
Consider the addition of two arbitrary independent angular momenta $\op{\vect J}_1$ and $\op{\vect J}_2$:
\[\op{\vect J}=\op{\vect J}_1+\op{\vect J}_2.\]
They satisfy
\[[\op J_{1i},\op J_{1j}]=\ii\hbar\epsilon_{ijk}\op J_{1k},\qquad
[\op J_{2i},\op J_{2j}]=\ii\hbar\epsilon_{ijk}\op J_{2k},\qquad
[\op J_{1i},\op J_{2j}]=0.\]
The respective eigenvalues are
\[J_1^2=j_1(j_1+1)\hbar^2,\quad J_{1z}=m_1\hbar,\quad m_1=-j_1,\ldots,j_1,\]
\[J_2^2=j_2(j_2+1)\hbar^2,\quad J_{2z}=m_2\hbar,\quad m_2=-j_2,\ldots,j_2.\]
It is easy to prove
\[[\op J_i,\op J_j]=\ii\hbar\epsilon_{ijk}\op J_k.\]
Thus $\op{\vect J}$ is again an angular-momentum operator, called the total angular momentum, with eigenvalues
\[J^2=j(j+1)\hbar^2,\qquad J_z=m\hbar,\qquad m=-j,\ldots,j.\]
The problem is to find the relation between $\ket{j,m}$ and $\ket{j_1,m_1}\ket{j_2,m_2}$.

\subsection*{1) Two representations}
\textbf{Uncoupled representation.} Since the two angular momenta commute with one another, the four operators $\op{\vect J}_1^{\,2}$, $\op J_{1z}$, $\op{\vect J}_2^{\,2}$, $\op J_{2z}$ have common eigenvectors $\ket{j_1m_1j_2m_2}$, which form the uncoupled representation:
\[\begin{aligned}
\op{\vect J}_1^{\,2}\ket{j_1m_1j_2m_2}&=j_1(j_1+1)\hbar^2\ket{j_1m_1j_2m_2},\\
\op J_{1z}\ket{j_1m_1j_2m_2}&=m_1\hbar\ket{j_1m_1j_2m_2},\\
\op{\vect J}_2^{\,2}\ket{j_1m_1j_2m_2}&=j_2(j_2+1)\hbar^2\ket{j_1m_1j_2m_2},\\
\op J_{2z}\ket{j_1m_1j_2m_2}&=m_2\hbar\ket{j_1m_1j_2m_2}.
\end{aligned}\]
But $[\op{\vect J}^{\,2},\op J_{1z}]\ne0$ and $[\op{\vect J}^{\,2},\op J_{2z}]\ne0$. Consequently, $\op{\vect J}^{\,2}$ is not diagonal in this basis, making its eigenvalues inconvenient to obtain directly.

\textbf{Coupled representation.} Since
\[[\op{\vect J}^{\,2},\op{\vect J}_1^{\,2}]
=[\op{\vect J}^{\,2},\op{\vect J}_2^{\,2}]
=[\op{\vect J}^{\,2},\op J_z]=0,\]
and the remaining relevant pairs also commute, the four operators $\op{\vect J}_1^{\,2}$, $\op{\vect J}_2^{\,2}$, $\op{\vect J}^{\,2}$, and $\op J_z$ have common eigenvectors $\ket{j_1j_2jm}$, which form the coupled representation:
\[\begin{aligned}
\op{\vect J}_1^{\,2}\ket{j_1j_2jm}&=j_1(j_1+1)\hbar^2\ket{j_1j_2jm},\\
\op{\vect J}_2^{\,2}\ket{j_1j_2jm}&=j_2(j_2+1)\hbar^2\ket{j_1j_2jm},\\
\op{\vect J}^{\,2}\ket{j_1j_2jm}&=j(j+1)\hbar^2\ket{j_1j_2jm},\\
\op J_z\ket{j_1j_2jm}&=m\hbar\ket{j_1j_2jm}.
\end{aligned}\]
Because the eigenvalues and eigenvectors of the four uncoupled operators are known, obtaining those of total angular momentum requires connecting the two representations by a change of basis.

From completeness of the uncoupled basis,
\[\sum_{m_1,m_2}\ket{j_1m_1j_2m_2}\bra{j_1m_1j_2m_2}=\mathbf1,\]
the transformation is
\[\boxed{\ket{j_1j_2jm}=\sum_{m_1,m_2}\ket{j_1m_1j_2m_2}
\braket{j_1m_1j_2m_2}{j_1j_2jm}.}\]
The overlaps $\braket{j_1m_1j_2m_2}{j_1j_2jm}$ are the Clebsch--Gordan (CG) coefficients.

\subsection*{2) Eigenvalues $(j,m)$ of the coupled representation}
Apply $\op J_z=\op J_{1z}+\op J_{2z}$ to the expansion:
\begin{align*}
m\hbar\ket{j_1j_2jm}
&=\sum_{m_1,m_2}\braket{j_1m_1j_2m_2}{j_1j_2jm}
(\op J_{1z}+\op J_{2z})\ket{j_1m_1j_2m_2}\\
&=\sum_{m_1,m_2}(m_1+m_2)\hbar\,
\braket{j_1m_1j_2m_2}{j_1j_2jm}\ket{j_1m_1j_2m_2}.
\end{align*}
Thus
\[\sum_{m_1,m_2}(m-m_1-m_2)\hbar\,
\braket{j_1m_1j_2m_2}{j_1j_2jm}\ket{j_1m_1j_2m_2}=0.\]
The uncoupled basis vectors are linearly independent. Therefore for each term either the CG coefficient is zero or $m=m_1+m_2$. We seek nonzero CG coefficients, and hence
\[\boxed{m=m_1+m_2.}\]
Next consider the possible values of $j$. Let $j_{\min}\le j\le j_{\max}$. The maximum is
\[j_{\max}=m_{\max}=(m_1)_{\max}+(m_2)_{\max}=j_1+j_2.\]
The dimension of the uncoupled representation is
\[D=(2j_1+1)(2j_2+1),\]
while the dimension of the coupled representation is
\[D=\sum_{j=j_{\min}}^{j_{\max}}(2j+1)
=(j_{\max}+1)^2-j_{\min}^2.\]
Equating the dimensions,
\[(2j_1+1)(2j_2+1)=(j_1+j_2+1)^2-j_{\min}^2,\]
gives $j_{\min}^2=(j_1-j_2)^2$ and $j_{\min}=|j_1-j_2|$.
Thus for fixed $(j_1,j_2)$,
\[\boxed{\begin{gathered}
J^2=j(j+1)\hbar^2,\qquad j=|j_1-j_2|,|j_1-j_2|+1,\ldots,j_1+j_2,\\
J_z=m\hbar,\qquad m=m_1+m_2.
\end{gathered}}\]

\subsection*{3) Eigenstates $\ket{j_1j_2jm}$ of the coupled representation}
The central task is to calculate the CG coefficients. We do not undertake a general treatment; standard tables are available.

\textbf{Example 1: spin and orbital angular momentum.} Let $\op{\vect J}_1=\op{\vect L}$, $\op{\vect J}_2=\op{\vect S}$, and $\op{\vect J}=\op{\vect L}+\op{\vect S}$. The coupled basis is $\ket{l,\tfrac12;j,m}$ and the uncoupled basis is $\ket{l,m_l;\tfrac12,m_s}$.
The transformation reads
\[\ket{l,\tfrac12;j,m}
=\sum_{m_l,m_s}C_{m_lm_s}\ket{l,m_l;\tfrac12,m_s}
=\sum_{m_l}\big[A_{m_l}\ket{l,m_l;\tfrac12,\tfrac12}
+B_{m_l}\ket{l,m_l;\tfrac12,-\tfrac12}\big].\]
Because $m=m_l+m_s$,
\[m_l=m-m_s=\begin{cases}
m-\frac12,&m_s=\frac12,\\
m+\frac12,&m_s=-\frac12,
\end{cases}\]
there are only two terms:
\[\boxed{\ket{l,\tfrac12;j,m}
=A\ket{l,m-\tfrac12;\tfrac12,\tfrac12}
+B\ket{l,m+\tfrac12;\tfrac12,-\tfrac12}.}\]
We determine the coefficients $A,B$.

In the spin subspace of the uncoupled representation,
\begin{align*}
\op{\vect J}^{\,2}&=\op{\vect L}^{\,2}+\op{\vect S}^{\,2}
+2\op{\vect L}\cdot\op{\vect S}\\
&=\begin{pmatrix}
\op{\vect L}^{\,2}+\frac34\hbar^2+\hbar\op L_z&\hbar\op L_-\\
\hbar\op L_+&\op{\vect L}^{\,2}+\frac34\hbar^2-\hbar\op L_z
\end{pmatrix}.
\end{align*}
Here $\op L_\pm=\op L_x\pm\ii\op L_y$ are the orbital raising and lowering operators. The uncoupled states are represented as
\[\ket{l,m-\tfrac12;\tfrac12,\tfrac12}
=\begin{pmatrix}\ket{l,m-\tfrac12}\\0\end{pmatrix},\qquad
\ket{l,m+\tfrac12;\tfrac12,-\tfrac12}
=\begin{pmatrix}0\\\ket{l,m+\tfrac12}\end{pmatrix}.\]
Insert the expansion into the eigenvalue equation
\[\op{\vect J}^{\,2}\ket{l,\tfrac12;j,m}
=j(j+1)\hbar^2\ket{l,\tfrac12;j,m}.\]
It becomes
\begingroup\small
\begin{align*}
j(j+1)\hbar^2\big(A\ket{l,m-\tfrac12;\tfrac12,\tfrac12}
+B\ket{l,m+\tfrac12;\tfrac12,-\tfrac12}\big)
={}&(\op{\vect L}^{\,2}+\op{\vect S}^{\,2}+2\op{\vect L}\cdot\op{\vect S})\\[-1ex]
&\quad\times\big(A\ket{l,m-\tfrac12;\tfrac12,\tfrac12}
+B\ket{l,m+\tfrac12;\tfrac12,-\tfrac12}\big).
\end{align*}
\endgroup
In the spin subspace this gives
\[\begin{cases}
A\left(\op{\vect L}^{\,2}+\frac34\hbar^2+\hbar\op L_z-j(j+1)\hbar^2\right)\ket{l,m-\tfrac12}
+B\hbar\op L_-\ket{l,m+\tfrac12}=0,\\[1ex]
A\hbar\op L_+\ket{l,m-\tfrac12}
+B\left(\op{\vect L}^{\,2}+\frac34\hbar^2-\hbar\op L_z-j(j+1)\hbar^2\right)\ket{l,m+\tfrac12}=0.
\end{cases}\]
Using the orbital-angular-momentum eigenvalues and ladder-operator coefficients yields
\[\begin{cases}
\displaystyle\left[l(l+1)+\frac34+\left(m-\frac12\right)-j(j+1)\right]A
+\sqrt{\left(l-m+\frac12\right)\left(l+m+\frac12\right)}\,B=0,\\[1.1ex]
\displaystyle\sqrt{\left(l+m+\frac12\right)\left(l-m+\frac12\right)}\,A
+\left[l(l+1)+\frac34-\left(m+\frac12\right)-j(j+1)\right]B=0.
\end{cases}\]
These equations determine $A,B$. The allowed total angular momenta are $j=l-\frac12$ and $j=l+\frac12$ (where permitted).

For $j=l+\frac12$, normalization yields
\[\boxed{\ket{l,\tfrac12;j,m}
=\sqrt{\frac{l+m+\frac12}{2l+1}}\ket{l,m-\tfrac12;\tfrac12,\tfrac12}
+\sqrt{\frac{l-m+\frac12}{2l+1}}\ket{l,m+\tfrac12;\tfrac12,-\tfrac12}.}\]
For $j=l-\frac12$,
\[\boxed{\ket{l,\tfrac12;j,m}
=-\sqrt{\frac{l-m+\frac12}{2l+1}}\ket{l,m-\tfrac12;\tfrac12,\tfrac12}
+\sqrt{\frac{l+m+\frac12}{2l+1}}\ket{l,m+\tfrac12;\tfrac12,-\tfrac12}.}\]

\textbf{Example 2: coupling two electron spins.} Let
\[\op{\vect s}=\op{\vect s}_1+\op{\vect s}_2,\qquad
s_1=s_2=\frac12,\qquad m_1,m_2=\pm\frac12,\qquad s=0,1.\]
The uncoupled basis states are
\[\begin{gathered}
\ket{\tfrac12,\tfrac12}\ket{\tfrac12,\tfrac12},\quad
\ket{\tfrac12,\tfrac12}\ket{\tfrac12,-\tfrac12},\\
\ket{\tfrac12,-\tfrac12}\ket{\tfrac12,\tfrac12},\quad
\ket{\tfrac12,-\tfrac12}\ket{\tfrac12,-\tfrac12}.
\end{gathered}\]
The coupled basis is $\ket{\tfrac12,\tfrac12;s,m}$ with $(s,m)=(1,1),(1,0),(1,-1),(0,0)$. As in Example 1, coupled states can be expressed in the uncoupled basis.

For $s=1$, the \emph{spin triplet} is
\begin{align*}
\ket{\tfrac12,\tfrac12;1,1}
&=\ket{\tfrac12,\tfrac12}\ket{\tfrac12,\tfrac12},\\
\ket{\tfrac12,\tfrac12;1,0}
&=\frac1{\sqrt2}\left(\ket{\tfrac12,\tfrac12}\ket{\tfrac12,-\tfrac12}
+\ket{\tfrac12,-\tfrac12}\ket{\tfrac12,\tfrac12}\right),\\
\ket{\tfrac12,\tfrac12;1,-1}
&=\ket{\tfrac12,-\tfrac12}\ket{\tfrac12,-\tfrac12}.
\end{align*}
For $s=0$, the \emph{spin singlet} is
\[\boxed{\ket{\tfrac12,\tfrac12;0,0}
=\frac1{\sqrt2}\left(\ket{\tfrac12,\tfrac12}\ket{\tfrac12,-\tfrac12}
-\ket{\tfrac12,-\tfrac12}\ket{\tfrac12,\tfrac12}\right).}\]